\paragraph{Native nucleon weak form factors.}
The literal native quark weak generator, contracted with the internal
Pauli currents, acts on the complete three-quark colour sector. Its
spin-half/isospin-half nucleon quartet has exact vector normalization
$g_V=1$. An integer epsilon-colour vector gives the co-located axial
coefficient $g_A=5/3$ and squared excited-state leakage $8/9$.
For the calibrated source, the separated-colour complement is above the
local trial by
\[
 \gamma=\frac{8U}{3}-\frac{4A}{9}
       -\frac34(J_s+J_w)-3t>0.
\]
Its mixing amplitude is at most $3t/\gamma$. Since
$\|I_-J_+^z\|\leq6$, the complete ground-state coefficient satisfies
$|g_A-5/3|\leq36t/\gamma$, giving
$g_A\in[1.665734201267,1.667599132066]$.
The PDG neutron-decay magnitude $1.2753\pm0.0013$\footnote{Particle Data
Group, neutron listing, 2025 update, p.~11:
\url{https://pdg.lbl.gov/2025/listings/rpp2025-list-n.pdf}.}
is excluded by this
interval. That held-out comparison rejects this bare-current physical
identification despite the fitted hadron masses. The nonzero excited
current channels are retained in the subsequent conversion calculation.
The exact witness and bound are in
\path{model/openwave_baryon_weak_formfactors}; no renormalized QCD axial
current or physical weak lifetime is inferred.
The local colour-singlet space decomposes as
$\operatorname{Sym}^3(\mathbb C^2\otimes\mathbb C^2)
=(\tfrac12,\tfrac12)\oplus(\tfrac32,\tfrac32)$, each with multiplicity
one. Schur's lemma therefore excludes changing this local axial
coefficient by a scalar deformation preserving spin and weak symmetry.
Additional state mixing or a justified effective-current correction is
required to change that result.
